% Resposta corrigida: Problema 2.20 — Fluxo em tubo rugoso (Buckingham Pi)
\subsection*{Identificação das Variáveis e Dimensões}

O problema postula que a queda de pressão $p$ ao longo do tubo depende das seguintes variáveis:
\[
p = p(v,\,\rho,\,\mu,\,l,\,d,\,e)
\]
onde $v$ é a velocidade do fluido, $\rho$ a densidade, $\mu$ a viscosidade dinâmica, $l$ o comprimento do tubo, $d$ o diâmetro e $e$ a rugosidade média.

Com $n = 7$ variáveis e $m = 3$ dimensões fundamentais ($M$, $L$, $T$), o teorema de Buckingham Pi fornece:
\[
k = n - m = 7 - 3 = 4 \text{ grupos adimensionais.}
\]

As dimensões das variáveis são:
\begin{table}[h!]
\centering
\begin{tabular}{lcc}
\toprule
\textbf{Variável} & \textbf{Símbolo} & \textbf{Dimensão} \\\midrule
Queda de pressão & $p$ & $M L^{-1} T^{-2}$ \\
Velocidade do fluido & $v$ & $L T^{-1}$ \\
Densidade & $\rho$ & $M L^{-3}$ \\
Viscosidade dinâmica & $\mu$ & $M L^{-1} T^{-1}$ \\
Comprimento do tubo & $l$ & $L$ \\
Diâmetro do tubo & $d$ & $L$ \\
Rugosidade média & $e$ & $L$ \\\bottomrule
\end{tabular}
\end{table}

\subsection*{Aplicação do Teorema de Buckingham Pi}

Escolhem-se $\rho$, $v$ e $d$ como variáveis de repetição (contêm juntas $M$, $L$ e $T$ sem dependência linear). Os quatro grupos Pi são construídos adicionando cada variável restante:

\textbf{Grupo $\Pi_1$} (com $\mu$):
\[
\Pi_1 = \rho^a\, v^b\, d^c\, \mu
\]
Impondo adimensionalidade ($M^0 L^0 T^0$):
\[
(M L^{-3})^a (L T^{-1})^b (L)^c (M L^{-1} T^{-1}) = M^0 L^0 T^0
\]
\[
\begin{cases} M:\; a + 1 = 0 \implies a = -1 \\ T:\; -b - 1 = 0 \implies b = -1 \\ L:\; -3a + b + c - 1 = 0 \implies c = -1 \end{cases}
\]
\[
\Pi_1 = \frac{\mu}{\rho v d} = \frac{1}{Re}
\]
onde $Re = \rho v d / \mu$ é o número de Reynolds.

\textbf{Grupo $\Pi_2$} (com $l$):
\[
\Pi_2 = \frac{l}{d}
\]
(razão comprimento/diâmetro, puramente geométrica).

\textbf{Grupo $\Pi_3$} (com $e$):
\[
\Pi_3 = \frac{e}{d}
\]
(rugosidade relativa).

\textbf{Grupo $\Pi_4$} (com $p$):
\[
\Pi_4 = \rho^a\, v^b\, d^c\, p
\]
\[
(M L^{-3})^a (L T^{-1})^b (L)^c (M L^{-1} T^{-2}) = M^0 L^0 T^0
\]
\[
\begin{cases} M:\; a + 1 = 0 \implies a = -1 \\ T:\; -b - 2 = 0 \implies b = -2 \\ L:\; -3a + b + c - 1 = 0 \implies c = 0 \end{cases}
\]
\[
\Pi_4 = \frac{p}{\rho v^2}
\]

\subsection*{Resultado}

A relação adimensional geral é:
\[
\frac{p}{\rho v^2} = \Phi\!\left(Re,\; \frac{l}{d},\; \frac{e}{d}\right)
\]

Este resultado é consistente com a equação de Darcy-Weisbach, que mostra que a queda de pressão em tubos rugosos depende do número de Reynolds, da razão comprimento/diâmetro e da rugosidade relativa. A análise dimensional reduz um problema com 7 variáveis a uma relação entre 4 grupos adimensionais.

A Segunda Lei de Newton estabelece a seguinte Equação Diferencial Ordinária (EDO):
\begin{equation}
    m \frac{dv}{dt} = mg - cv^2
\end{equation}

Definimos a velocidade terminal ($v_t$) como o estado onde a aceleração é nula ($mg = cv_t^2$). Assim, podemos expressar a constante de arrasto como $c = \frac{mg}{v_t^2}$. Substituindo na EDO original, obtemos a forma normalizada:
\begin{equation}
    \frac{dv}{dt} = g \left( 1 - \frac{v^2}{v_t^2} \right)
\end{equation}

\subsection*{Separação de Variáveis e Integração}
Para determinar o tempo $t$ necessário para atingir uma velocidade $v$, isolamos os termos e integramos nos limites definidos pela condição inicial:
\begin{equation}
    \int_{v_0}^{v} \frac{dv}{1 - \left( \frac{v}{v_t} \right)^2} = \int_{0}^{t} g \, dt
\end{equation}

Utilizando a expansão por frações parciais para o termo à esquerda:
\begin{equation}
    \frac{1}{1 - (v/v_t)^2} = \frac{v_t^2}{v_t^2 - v^2} = \frac{v_t}{2} \left( \frac{1}{v_t + v} + \frac{1}{v_t - v} \right)
\end{equation}

Integrando ambos os lados:
\begin{equation}
    \frac{v_t}{2} \left[ \ln \left( \frac{v_t + v}{v_t - v} \right) \right]_{v_0}^{v} = gt
\end{equation}

\subsection*{Solução Analítica}
Aplicando os limites de integração superiores e inferiores, obtemos a expressão para o tempo $t$:
\begin{equation}
    t = \frac{v_t}{2g} \left[ \ln \left( \frac{v_t + v}{v_t - v} \right) - \ln \left( \frac{v_t + v_0}{v_t - v_0} \right) \right]
\end{equation}

Utilizando as propriedades de logaritmo ($\ln A - \ln B = \ln \frac{A}{B}$), a solução final compacta é:
\begin{equation}
    t = \frac{v_t}{2g} \ln \left( \frac{(v_t + v)(v_t - v_0)}{(v_t - v)(v_t + v_0)} \right)
\end{equation}

\subsection*{Análise de Convergência}
Seja $\alpha$ a fração da velocidade terminal que se deseja atingir ($v = \alpha v_t$, com $0 < \alpha < 1$). O tempo $t_\alpha$ é dado por:
\begin{equation}
    t_{\alpha} = \frac{v_t}{2g} \ln \left( \frac{(1 + \alpha)(v_t - v_0)}{(1 - \alpha)(v_t + v_0)} \right)
\end{equation}

\textbf{Conclusão:} A presença da velocidade inicial $v_0$ atua como uma translação na escala temporal. Se $v_0 = 0$, a solução converge para o modelo de queda livre padrão. Se $v_0 > v_t$, o objeto desacelerará até atingir a velocidade terminal, mantendo a validade física da solução logarítmica (em módulo).
